\documentclass[10pt]{article}

\usepackage[margin=0.78in]{geometry}
\usepackage{amsmath,amssymb,amsfonts}
\usepackage{booktabs}
\usepackage{microtype}
\usepackage[hidelinks]{hyperref}
\usepackage{enumitem}
\setlist{nosep,leftmargin=*}

\title{Inductive Algebraic Resonance Memory: A Non-Attention Causal Language Architecture}
\author{Godson Johnson}
\date{}

\newcommand{\R}{\mathbb{R}}
\newcommand{\od}{\odot}
\newcommand{\oslashop}{\oslash}

\begin{document}
\maketitle

\begin{abstract}
Inductive Algebraic Resonance Memory (IARM) is a causal language-model architecture that replaces pairwise self-attention with learned low-rank query operators and a prefix-normalized memory scan. The architecture preserves the residual, normalized, feed-forward structure of modern decoder language models, but its sequence mixer contains no key projection, no query-key score matrix, and no softmax attention matrix. Each layer combines depthwise causal local convolution, rotary query phases, gated low-rank query transformations, positive feature construction, cumulative feature-value memory, output gating, and a SwiGLU feed-forward block. This note gives the mathematical specification of the IARM language model and clarifies its relation to attention-based and linear-time sequence models.
\end{abstract}

\section{Motivation}
Autoregressive language models commonly use causal self-attention, where each position forms query-key scores against all earlier positions and then mixes values through a normalized attention matrix \cite{vaswani2017attention}. This mechanism is expressive, but it explicitly constructs token-token interactions. Several alternatives have sought linear-time sequence mixing, including kernelized attention \cite{katharopoulos2020transformers}, random-feature attention \cite{choromanski2021rethinking}, recurrent state-space models \cite{gu2023mamba}, and convolutional or gated convolutional sequence models \cite{chollet2017xception,dauphin2017language}.

IARM follows a different route. It does not approximate causal self-attention. Instead, it constructs a learned causal memory from positive transformed query features and values. The current token representation contributes to a cumulative memory state, and the readout is a coordinatewise normalized prefix statistic. The central architectural idea is that sequence context can be represented by an evolving algebraic memory rather than by an explicit attention matrix.

\section{Model Overview}
Let the token sequence be $x_{1:T}$ with vocabulary size $V$. The model embeds tokens into $d$ dimensions and applies dropout,
\begin{equation}
H^{(0)} = \operatorname{Dropout}(E[x_{1:T}]), \qquad E \in \R^{V\times d}. 
\end{equation}
The model consists of $L$ identical residual blocks. For layer $\ell$, the update is
\begin{align}
\widetilde H^{(\ell)} &= \operatorname{IARMUnit}\!\left(H^{(\ell-1)}\right), \\
H^{(\ell)} &= \widetilde H^{(\ell)} + \operatorname{SwiGLU}\!\left(\operatorname{RMSNorm}(\widetilde H^{(\ell)})\right).
\end{align}
After the final block, a root-mean-square normalization \cite{zhang2019root} is applied and logits are produced with a tied language-model head \cite{inan2017tying,press2017using},
\begin{equation}
Z = \operatorname{RMSNorm}(H^{(L)}), \qquad \operatorname{logits}_t = E Z_t.
\end{equation}
The feed-forward sublayer uses a SwiGLU nonlinearity, following the gated feed-forward design used in modern language models \cite{shazeer2020glu}.

\section{Inductive Algebraic Resonance Unit}
For a layer input $X\in\R^{B\times T\times d}$, the unit first computes
\begin{equation}
\bar X = \operatorname{RMSNorm}(X).
\end{equation}
A depthwise causal convolution gives a local temporal signal
\begin{equation}
L = \operatorname{DepthwiseCausalConv}(\bar X),
\end{equation}
where the convolution is channelwise, uses kernel width $k$, and truncates the right padding so that position $t$ depends only on positions at or before $t$.

The normalized states are projected into queries and values,
\begin{equation}
Q = W_q \bar X, \qquad V = W_v \bar X,
\end{equation}
then reshaped into $A$ heads of width $d_h=d/A$. Rotary positional embeddings are applied to queries only \cite{su2021roformer},
\begin{equation}
Q \leftarrow \operatorname{RoPE}(Q).
\end{equation}
There is no key tensor. Consequently, the unit has no query-key dot product and no attention matrix.

\subsection{Gated Low-Rank Query Operators}
For each head $a\in\{1,\ldots,A\}$ and operator index $o\in\{1,\ldots,O\}$, IARM learns two low-rank factors
\begin{equation}
U_{a,o},R_{a,o}\in\R^{d_h\times r},
\end{equation}
which define an implicit rank-$r$ operator $U_{a,o}R_{a,o}^{\top}$. A gate over operators is predicted from the normalized hidden state,
\begin{equation}
G_{b,a,t,o}=
\frac{\exp((W_g\bar X_{b,t})_{a,o})}
{\sum_{o'=1}^{O}\exp((W_g\bar X_{b,t})_{a,o'})}.
\end{equation}
The query displacement is a gated mixture of low-rank operator actions,
\begin{equation}
\Delta Q_{b,a,t}
=
\sum_{o=1}^{O}
G_{b,a,t,o}\,U_{a,o}\left(R_{a,o}^{\top}Q_{b,a,t}\right).
\end{equation}
The transformed query is
\begin{equation}
\widehat Q_{b,a,t}=Q_{b,a,t}+\frac{\Delta Q_{b,a,t}}{\sqrt r}.
\end{equation}
This operator pathway is related in spirit to low-rank adaptation and low-rank parameterization \cite{hu2022lora}, but here the low-rank factors are intrinsic sequence-mixing operators rather than adapters added to an existing attention model.

\subsection{Positive Features and Causal Memory Scan}
The transformed query is mapped to a positive feature vector using the shifted ELU feature map \cite{clevert2016fast},
\begin{equation}
\Phi_{b,a,t}=\operatorname{ELU}(\widehat Q_{b,a,t})+1.
\end{equation}
The causal memory is then built through cumulative coordinatewise sums,
\begin{align}
N_{b,a,t} &= \sum_{s\leq t}\Phi_{b,a,s}\od V_{b,a,s}, \\
D_{b,a,t} &= \max\left(\sum_{s\leq t}\Phi_{b,a,s},\varepsilon\right),
\end{align}
where the maximum is applied coordinatewise for numerical stability. The memory read is
\begin{equation}
M_{b,a,t}=N_{b,a,t}\oslashop D_{b,a,t}.
\end{equation}
This operation is causal because all sums are restricted to $s\leq t$. It is not attention: the readout is not a weighted sum over values using scores from the current position to every past position. It is a prefix-normalized feature-value statistic.

\subsection{Output Gating and Residual Return}
The multi-head memory output is reshaped to $\R^d$ and combined with the local convolution signal. A learned output gate controls the emitted sequence-mixing signal,
\begin{align}
\Gamma &= \sigma(W_o\bar X),\\
Y &= W_{\mathrm{out}}\left(\Gamma\od (\operatorname{MergeHeads}(M)+L)\right).
\end{align}
The unit returns a residual update,
\begin{equation}
\operatorname{IARMUnit}(X)=X+\operatorname{Dropout}(Y).
\end{equation}

\section{Distinction from Causal Self-Attention}
Causal self-attention computes
\begin{equation}
\operatorname{Attn}(Q,K,V)_t
=
\sum_{s\leq t}
\operatorname{softmax}_{s}\left(\frac{Q_tK_s^{\top}}{\sqrt{d_h}}\right)V_s.
\end{equation}
IARM does not compute this expression or an approximation of it. Its memory read is instead
\begin{equation}
M_t=
\frac{\sum_{s\leq t}\Phi_s\od V_s}
{\max\left(\sum_{s\leq t}\Phi_s,\varepsilon\right)}.
\end{equation}
The distinction is structural. Attention uses current-position queries to score all past keys. IARM accumulates transformed positive query features and values into a causal memory. It therefore has no explicit token-token score matrix and no softmax over previous positions.

\section{Complexity and Parameterization}
For sequence length $T$, width $d$, number of heads $A$, operator count $O$, and rank $r$, the scan part of IARM is linear in $T$. The operator transformation contributes approximately $O(BTAOrd_h)$ operations, while the cumulative memory scan contributes $O(BTAd_h)$. Since $d=Ad_h$, the dominant sequence dependence is linear in $T$ for fixed $O$ and $r$. This contrasts with the $O(T^2)$ score matrix used in standard full causal attention.

The unit also includes a depthwise convolution with cost $O(BTdk)$ and a SwiGLU block with the usual feed-forward cost determined by the hidden dimension. The architecture therefore separates short-range local filtering, algebraic operator transformation, and prefix memory accumulation.

\section{Discussion}
IARM is best understood as a non-attention causal language architecture. Its inductive bias is not pairwise retrieval but progressive memory formation: each position contributes a positive transformed query feature and a value to a cumulative memory; the read is a normalized prefix statistic. The low-rank operator bank gives the model a learned algebraic mechanism for reshaping query geometry before memory formation. The local convolution supplies immediate neighborhood structure, while the SwiGLU block supplies channelwise nonlinear transformation.

This design keeps several stabilizing ingredients from modern decoder models, including residual connections, RMS normalization, rotary phases, gated feed-forward layers, dropout, and tied embeddings. Its defining difference is the absence of causal self-attention. The model's sequence mixer is a causal resonance-memory scan built from learned operator-transformed features.

\begin{thebibliography}{99}

\bibitem{vaswani2017attention}
A. Vaswani, N. Shazeer, N. Parmar, J. Uszkoreit, L. Jones, A. N. Gomez, L. Kaiser, and I. Polosukhin.
Attention is all you need.
\emph{Advances in Neural Information Processing Systems}, 2017.

\bibitem{katharopoulos2020transformers}
A. Katharopoulos, A. Vyas, N. Pappas, and F. Fleuret.
Transformers are RNNs: Fast autoregressive transformers with linear attention.
\emph{Proceedings of ICML}, 2020.

\bibitem{choromanski2021rethinking}
K. Choromanski, V. Likhosherstov, D. Dohan, X. Song, A. Gane, T. Sarlos, P. Hawkins, J. Davis, A. Mohiuddin, L. Kaiser, D. Belanger, L. Colwell, and A. Weller.
Rethinking attention with performers.
\emph{Proceedings of ICLR}, 2021.

\bibitem{gu2023mamba}
A. Gu and T. Dao.
Mamba: Linear-time sequence modeling with selective state spaces.
\emph{arXiv:2312.00752}, 2023.

\bibitem{zhang2019root}
B. Zhang and R. Sennrich.
Root mean square layer normalization.
\emph{Advances in Neural Information Processing Systems}, 2019.

\bibitem{su2021roformer}
J. Su, Y. Lu, S. Pan, B. Wen, and Y. Liu.
RoFormer: Enhanced transformer with rotary position embedding.
\emph{arXiv:2104.09864}, 2021.

\bibitem{shazeer2020glu}
N. Shazeer.
GLU variants improve transformer.
\emph{arXiv:2002.05202}, 2020.

\bibitem{chollet2017xception}
F. Chollet.
Xception: Deep learning with depthwise separable convolutions.
\emph{Proceedings of CVPR}, 2017.

\bibitem{dauphin2017language}
Y. N. Dauphin, A. Fan, M. Auli, and D. Grangier.
Language modeling with gated convolutional networks.
\emph{Proceedings of ICML}, 2017.

\bibitem{clevert2016fast}
D. A. Clevert, T. Unterthiner, and S. Hochreiter.
Fast and accurate deep network learning by exponential linear units (ELUs).
\emph{Proceedings of ICLR}, 2016.

\bibitem{inan2017tying}
H. Inan, K. Khosravi, and R. Socher.
Tying word vectors and word classifiers: A loss framework for language modeling.
\emph{Proceedings of ICLR}, 2017.

\bibitem{press2017using}
O. Press and L. Wolf.
Using the output embedding to improve language models.
\emph{Proceedings of EACL}, 2017.

\bibitem{hu2022lora}
E. J. Hu, Y. Shen, P. Wallis, Z. Allen-Zhu, Y. Li, S. Wang, L. Wang, and W. Chen.
LoRA: Low-rank adaptation of large language models.
\emph{Proceedings of ICLR}, 2022.

\end{thebibliography}

\end{document}
